\hypertarget{discussion-and-limitations}{%
\section{6. Discussion and
limitations}\label{discussion-and-limitations}}

The measured approximation follows the reference with a modest final gap
on this workload, halves matrix-state storage, and lowers measured step
time. That result is useful as an implementation operating point. It
does not show that warm tracking, BF16 storage, or the particular repair
schedule is generally optimal. Storage and compute changes also alter
numerical trajectories.

The evidence covers one model scale, one dataset cohort/order, two model
initializations, one GPU family, and a validation panel used during
development. There are no generalization claims across architectures,
corpora, or hardware, and no untouched test result or statistical
equivalence test. Short QR20/QR50 controls cannot determine a
final-horizon ranking. Synthetic QR inputs and one-factor probes do not
represent the trained distribution of all factor pairs.

The primary time comparison uses the original audited FP32 QR10
implementation. Our isolated reuse profile narrows one implementation
attribution question but does not replace a newly trained
efficient-reference comparison. There is no direct matched evaluation
against subspace or adaptive-refresh methods, and the warm step has no
standalone diagonalization guarantee.

The early-gap pattern is qualitatively consistent with Eschenhagen et
al.'s finding that accurate bases mattered more early in their Imagewoof
ViT experiments (§4.4). That different workload does not establish the
mechanism of our gap. Our fixed schedule and reference-calibrated LR
leave that question open. Ignoring startup and other excluded work, the
seed205 mean-step ratio assigns 10,000 candidate calls the same measured
steady cost as approximately 8,134 stock calls. Whether a reference with
an appropriately shortened LR horizon would reach the same loss is
unmeasured. This arithmetic is not a time-to-quality result.
